% ============================================================
% CHAPTER 11 — Project 8: Flutter Setup and First Mobile App
% ============================================================
\chapter{Project 8: Flutter Setup and First Mobile App}\index{Flutter}\index{Dart}\index{Riverpod}\index{Dio}

\section*{What You'll Build}
A working Flutter mobile application that:
\begin{itemize}
  \item Runs on Android and iOS from a single codebase
  \item Has a login screen and a notes list (static UI)
  \item Uses the 0-code pattern with \texttt{\_CONTEXT.md} for Flutter/Dart
\end{itemize}

\textbf{Estimated time}: 60--90 minutes\\
\textbf{Prerequisite}: Working Notes API backend (Chapters~6--10)

\begin{notebox}
\textbf{Tooling Box --- Stack chosen for this example.}
\begin{itemize}
  \item \textbf{Framework:} Flutter~3.x (Dart~3.x)
  \item \textbf{State Management:} Riverpod~2.x
  \item \textbf{Routing:} GoRouter~14.x
  \item \textbf{HTTP Client:} Dio
\end{itemize}
\textbf{Equivalent alternatives:} React Native, Kotlin Multiplatform, .NET MAUI. The \textbf{pattern} (declarative UI, state management, navigation, API calls to the backend) is common to all cross-platform mobile frameworks. If you choose React Native, the ADLC method and \texttt{\_CONTEXT.md} work in exactly the same way.
\end{notebox}

% ---------------------------------------------------------
\section{Installing Flutter}

\begin{versionbox}
The instructions in this chapter are verified with Flutter~3.x, Dart~3.x, Riverpod~2.x, and GoRouter~14.x. The Flutter ecosystem evolves rapidly: packages change APIs between minor versions. If the AI flags a deprecated method or suggests a newer version of a package, trust the correction --- the architectural pattern stays the same. Always run \texttt{flutter pub get} after changes to \texttt{pubspec.yaml}.
\end{versionbox}

\begin{praticabox}[title={Hands-on --- Installation}]
\textbf{Windows:}
\begin{lstlisting}[language=bash]
winget install -e --id Google.Flutter
\end{lstlisting}

\textbf{macOS:}
\begin{lstlisting}[language=bash]
brew install --cask flutter
\end{lstlisting}

\textbf{Linux:}
\begin{lstlisting}[language=bash]
sudo snap install flutter --classic
\end{lstlisting}

\textbf{Verification:}
\begin{lstlisting}[language=bash]
flutter doctor
\end{lstlisting}
You should have: Flutter SDK \checkmark, Android toolchain \checkmark, Connected device \checkmark.
\end{praticabox}

\begin{tipbox}
For Android development without the full Android Studio, you can use only the Android SDK command-line tools. But for your first projects, Android Studio with the built-in emulator is the simplest choice.
\end{tipbox}

Install the \textbf{Flutter} (Dart-Code) extension for VS~Code --- it enables IntelliSense, running, debugging, hot reload, and the widget inspector.

% ---------------------------------------------------------
\section{Creating the Project}

\begin{praticabox}[title={Hands-on --- Project scaffold}]
\begin{lstlisting}[language=bash]
cd notes-fullstack
flutter create notes_mobile
cd notes_mobile
\end{lstlisting}
Generated structure:
\begin{lstlisting}[language={}]
notes_mobile/
+-- lib/
|   +-- main.dart          <- Entry point
+-- android/               <- Android configuration
+-- ios/                   <- iOS configuration
+-- test/                  <- Tests
+-- pubspec.yaml           <- Dependencies
+-- analysis_options.yaml  <- Lint rules
\end{lstlisting}
\end{praticabox}

\begin{praticabox}[title={Hands-on --- Verify it works}]
\begin{lstlisting}[language=bash]
flutter run
# Or for quick testing:
flutter run -d chrome
\end{lstlisting}
\end{praticabox}

\begin{tipbox}
\texttt{flutter run -d chrome} is the fastest way to iterate. You do not need an emulator for most UI work.
\end{tipbox}

% ---------------------------------------------------------
\section{The Context for Flutter}

\begin{lstlisting}[language={},title={\texttt{\_CONTEXT.md} for Notes Mobile}]
# Project: Notes Mobile (Flutter)

## Purpose
Flutter mobile app connected to the Notes API backend.

## Technologies
- Flutter 3.x, Dart
- State Management: Riverpod 2
- HTTP: Dio
- Routing: GoRouter
- Local storage: flutter_secure_storage
- UI: Material Design 3

## Architecture
lib/
  main.dart
  app.dart
  config/
    api_config.dart      <- Backend URL, timeouts
    theme.dart           <- Material Design 3 theme
  models/
    user.dart / note.dart / api_response.dart
  services/
    api_service.dart     <- Configured Dio client
    auth_service.dart    <- Login, logout, refresh
    notes_service.dart   <- Notes CRUD
  providers/
    auth_provider.dart   <- Authentication state
    notes_provider.dart  <- Notes state
  screens/
    login_screen.dart
    dashboard_screen.dart
    note_detail_screen.dart
    note_form_screen.dart
  widgets/
    note_card.dart / empty_state.dart / error_banner.dart

## Dart/Flutter Conventions
- Files: snake_case.dart
- Classes: PascalCase
- Variables/functions: camelCase
- Widgets: each widget in a separate file
- State: Riverpod (NOT setState for shared state)
- Models: immutable classes with factory fromJson/toJson
- Max 3 widget nesting levels -> extract

## Constraints
- DO NOT use SharedPreferences for JWT tokens
- DO NOT use setState for state shared across screens
- DO NOT hardcode URLs. Use api_config.dart
- Every screen: loading, error, empty, data
- Material Design 3 with useMaterial3: true
\end{lstlisting}

% ---------------------------------------------------------
\section{\texttt{SKILL.md} for Flutter}

\begin{lstlisting}[language={},title={\texttt{SKILL.md} --- Flutter Mobile Developer}]
# Flutter Mobile Developer Skill

## Widget Structure
class NoteCard extends ConsumerWidget {
  final Note note;
  final VoidCallback? onTap;
  const NoteCard({super.key, required this.note, this.onTap});

  @override
  Widget build(BuildContext context, WidgetRef ref) {
    return Card(
      child: ListTile(
        title: Text(note.title),
        subtitle: Text(note.content, maxLines: 2),
        onTap: onTap,
      ),
    );
  }
}

## Riverpod Pattern
final authProvider =
  StateNotifierProvider<AuthNotifier, AuthState>((ref) {
    return AuthNotifier(ref.read(authServiceProvider));
  });

## Material 3 Design
- ColorScheme.fromSeed(seedColor: Colors.indigo)
- Theme.of(context).colorScheme for colors
- Card with elevation: 0 and border for a modern style
- FilledButton primary, OutlinedButton secondary
\end{lstlisting}

% ---------------------------------------------------------
\section{Generating the Base Structure}

\begin{praticabox}[title={Hands-on --- Generate the app with Copilot}]
\begin{lstlisting}[language={}]
Read _CONTEXT.md and SKILL.md. Generate the base structure.

Order:
 1. Dependencies in pubspec.yaml: dio, riverpod,
    flutter_riverpod, go_router, flutter_secure_storage
 2. lib/config/api_config.dart
 3. lib/config/theme.dart (Material Design 3, seed: indigo)
 4. lib/models/ (Note, User, ApiResponse) immutable
 5. lib/services/api_service.dart (Dio + interceptor)
 6. lib/providers/auth_provider.dart
 7. lib/app.dart with MaterialApp.router
 8. lib/screens/login_screen.dart (static UI)
 9. lib/screens/dashboard_screen.dart (static UI, mock data)
10. GoRouter with redirect for route protection
11. lib/main.dart
\end{lstlisting}
After the AI updates \texttt{pubspec.yaml}:
\begin{lstlisting}[language=bash]
flutter pub get
\end{lstlisting}
\end{praticabox}

% ---------------------------------------------------------
\section{Verification and Iteration}

\begin{praticabox}[title={Hands-on --- Start the app}]
\begin{lstlisting}[language=bash]
flutter run -d chrome
\end{lstlisting}
\end{praticabox}

\begin{center}
\begin{tabularx}{\textwidth}{l X}
\toprule
\textbf{Feature} & \textbf{Test} \\
\midrule
App starts & No compilation errors \\
Login screen & Google/GitHub buttons visible \\
Theme & Indigo Material~3 colors \\
Router & Login $\leftrightarrow$ dashboard navigation \\
Redirect & Direct access to /dashboard $\rightarrow$ /login \\
\bottomrule
\end{tabularx}
\end{center}

\textbf{Hot Reload}: with the app running, edit the text of a button and save. The app updates \textbf{instantly} --- the fastest feedback loop in mobile development.

\begin{checkpointbox}
The app starts, shows the login screen with the Material~3 theme, and navigation works.
\end{checkpointbox}

% ---------------------------------------------------------
\section{Dart in 10 Minutes}

You do not need to know Dart to work in 0-code. But recognizing the patterns helps you review:

\begin{lstlisting}[language=Java,title={Key Dart concepts}]
// Variables: strong typing, null safety
String name = 'Notes App';
int? count;                       // Nullable
final items = <Note>[];           // Typed list

// Asynchronous functions
Future<List<Note>> fetchNotes() async {
  final response = await dio.get('/api/notes');
  return (response.data['data'] as List)
      .map((json) => Note.fromJson(json))
      .toList();
}

// Classes
class Note {
  final String id;
  final String title;
  const Note({required this.id, required this.title});
}

// Widget = function that returns UI
class MyButton extends StatelessWidget {
  final String label;
  final VoidCallback onPressed;
  const MyButton({super.key,
    required this.label, required this.onPressed});

  @override
  Widget build(BuildContext context) {
    return FilledButton(
      onPressed: onPressed, child: Text(label));
  }
}
\end{lstlisting}

\begin{tipbox}
Dart looks a lot like Java/TypeScript. If you do not know them, do not worry: the AI writes the code, and you verify that it does what you asked.
\end{tipbox}

% ---------------------------------------------------------
\section*{Summary}

\begin{center}
\begin{tabularx}{\textwidth}{l X}
\toprule
\textbf{Aspect} & \textbf{Detail} \\
\midrule
Framework & Flutter~3.x with Dart \\
State & Riverpod~2 \\
Routing & GoRouter with redirect \\
Design & Material Design~3 \\
Structure & Monorepo notes-fullstack/notes\_mobile/ \\
\bottomrule
\end{tabularx}
\end{center}
